\hypertarget{sec:corporate-residuals}{%
\subsection{Changed obligations and unfinished production}\label{sec:corporate-residuals}}

A company can change its restoration estimate before a machine wears out, promise to repair sold products, or end a period with unfinished goods. Each condition changes a different part of its accounts. The following cases extend the preceding operating methods through those transitions. They retain exact illustrative amounts and separate economic evidence from bookkeeping.

Restoration, production and uncertain-tax methods use the retained historical EU text; assurance and mixed contracts use the inspected AASB compilations for synthetic 2024/2025 Australian for-profit reporters. The AASB 137 PDF covers annual periods beginning from 1 January 2023 to before 1 July 2026; its website gives a different end-period summary. Both admit these particular historical calendar years. A real application still needs its own source and period admission. Amounts are pretax unless an explicit tax calculation follows.

\hypertarget{restoration-estimates-change-the-asset-and-the-obligation-differently}{%
\subsubsection{Restoration estimates change the asset and the obligation differently}\label{restoration-estimates-change-the-asset-and-the-obligation-differently}}

The restoration record retains the related component and liability, cost/revaluation policy, gross cost/accumulated depreciation, remaining life/residual, current cash-flow estimate and rate, prior estimate identity, same-asset revaluation surplus/prior profit deficit, hypothetical cost-model carrying amount, impairment evidence and whole-class valuation review. IFRIC~1.2 \citep[IFRIC~1.1--8]{eu20231803} requires the original liability to have been included in the asset's cost. A standalone litigation obligation is not admitted to this asset-adjustment method.

Separate passage of time from a new estimate. Under the cost model, a liability increase adds to PPE; a decrease reduces PPE only to zero, with excess in profit. An increase prompts assessment of impairment indications and an IAS 36 test when required. Depreciate the adjusted carrying amount less admitted residual over the remaining useful life. Discount unwind is finance expense, never capitalised under IFRIC 1.8. Once useful life is exhausted, later estimate changes go to profit under 7 regardless of asset model. Actual settlement reduces the obligation and cash; a difference from the last carrying amount requires a dated remeasurement.

With cash 100, PPE 80 and provision 20, opening equity is 160. An admitted new provision 50 adds 30 to PPE. Recoverable amount 90 requires impairment 20: PPE 90, cash 100, liability 50 and equity 140. Four years remaining and zero residual gives annual depreciation 22.5. The subsequent depreciation-only illustration assumes a separately supplied zero discount rate and no further liability change; PPE 67.5 and equity 117.5 follow. In another branch, PPE 30/provision 80 and a new provision 30 require Dr provision 50 / Cr PPE 30 / Cr profit 20. Negative PPE 20 would be wrong. Cash 100 then equals liability 30 plus equity 70.

Under the revaluation model, a liability increase first consumes a same-asset surplus through OCI, then affects profit; a decrease reverses a same-asset profit deficit and otherwise increases OCI. A decrease exceeding the hypothetical cost-model carrying amount has an immediate profit component. Determine any necessary class revaluation first so that these splits use the correct history. With PPE 150/provision 50/surplus 20, an increase 35 gives OCI loss 20 and profit loss 15; PPE stays 150 only because a supplied class valuation review supports it. Cash 100 then gives equity 165 and liability 85. Surplus on another asset cannot be borrowed. In a separate decrease case, PPE 25, hypothetical cost 40, earlier profit deficit 15 and provision 80 → 30 yield excess 10 plus deficit reversal 15 in profit and 25 in OCI. Cash 100+PPE 25=provision 30+equity 95. Using 25 as the cost floor would compute the wrong split. The prior deficit explains why 25 and 40 differ; no fresh impairment is inferred from that difference.

A full cash history opens cash 100/PPE 120/provision 20/equity 200. Two-year straight-line life and zero residual give depreciation 60 in year 1. At 10\%, unwind 2 raises the provision 22; new supported estimate 33 adds 11 to PPE. Close PPE 71/provision 33/equity 138, with pretax loss 62 and no cash change. Year 2 depreciation 71 and unwind 3.3 give loss 74.3; actual settlement 36.3 leaves cash/equity 63.7 and zero PPE/provision. Across both years depreciation 131 plus finance 5.3 equals cumulative loss 136.3. The total cash bridge is exact; IAS 7 classification of restoration settlement needs a separate fact-specific policy. The method therefore returns a cash-classification obligation rather than inventing a complete categorized cash-flow statement. A fully exhausted PPE 0 with provision 20 → 25 gives expense 5 even if a historic revaluation reserve exists.

\hypertarget{assurance-is-a-measured-obligation-alongside-the-product-sale}{%
\subsubsection{Assurance is a measured obligation alongside the product sale}\label{assurance-is-a-measured-obligation-alongside-the-product-sale}}

The assurance record stores covered sales/population, legal or constructive promise, coverage/task findings, claim history, remaining-outcome law, timing/discount, reimbursement contract and virtual-certainty evidence. AASB 15.B 28--B 33 determines whether assurance, service or an inseparable combined service is present. AASB~137.14/24 \citep[paras.~14,24,36--62]{aasb137dated} admission precedes measurement. For a large population 39 uses weighted outcomes; individual claim probability does not replace the population outflow judgment. Correlation matters for distributions even when marginal means are sufficient to add expected cost. A single lawsuit uses 40's different analysis.

Consider 100 products with admitted per-product costs 0,10,40 and probabilities 3/4,1/5,1/20. The population estimate weights each possible repair cost and then sums across the covered products:

\[100\left(\frac34\,0+\frac15\,10+\frac1{20}\,40\right)=400.\]

A low claim probability for one product therefore does not eliminate the population obligation. The entire population has a present assurance duty, probable outflow and reliable estimate; discount is immaterial over the short horizon. Opening cash 1000 and inventory 600 are funded by equity 1600. Cash sales 2000 and inventory expense 600 accompany warranty expense/provision 400. An independently enforceable quarter-cost reimbursement is virtually certain on settlement, so recognise a separate asset 100/benefit 100. Do not net it against the provision. Paragraph 54 permits net expense presentation.

Pay claims 100 and receive reimbursement 25. A new remaining-claims estimate 250 replaces the old 300, releasing 50; reduce the remaining reimbursement 75 to 62.5, expense 12.5. Year 1 cash 2925+reimbursement 62.5=provision 250+pretax equity 2737.5. Profit 1137.5 follows from 2000 − 600 − 400+100+50 − 12.5. Direct operating cash 1925 also equals profit 1137.5+inventory release 600+provision increase 250 − receivable increase 62.5. Year 2 settlement 250 and receipt 62.5 leave cash/equity 2737.5 with no further profit under the realised amounts. Without virtual certainty, no reimbursement asset is recognised; retain the separate uncertainty. Missing claim evidence does not permit zero provision or the printed 400 by default.

\hypertarget{a-mixed-amendment-needs-two-revenue-calculations}{%
\subsubsection{A mixed amendment needs two revenue calculations}\label{a-mixed-amendment-needs-two-revenue-calculations}}

The modification record retains the original allocation/performance, enforceable change, remaining distinctness findings, amended allocations, faithful measures of progress and new receipts. AASB~15.21(c) \citep[paras.~18--21,39--45,73--90]{aasb15dated} requires consistent application of 21(a)/(b) to the affected promises; it does not supply a universal price split. The following exact case deliberately removes discount-allocation ambiguity.

A contract includes construction 100 and four distinct inspection visits 5 each, at their evidenced standalone prices. Inspections concern another operating line and can occur before the new construction finishes. The original expected construction cost 80; actual eligible cost 40 and faithful cost-to-cost progress give 50\% and revenue 50. One visit adds revenue 5 and cash expense 2. The 120 price was prepaid. The approved amendment adds integrated construction 40 and two visits 10; supported standalone allocations therefore become 140 and 30. New total construction cost 100 makes cumulative progress 40\%. The amended construction allocation gives
\[\text{cumulative revenue}=140\frac{40}{100}=56,\qquad\text{catch-up}=56-50=6.\]
This corrects only the construction promise. The first visit's 5 remains; five remaining distinct visits share 25 prospectively. The new 50 receipt itself earns no revenue.

With opening cash/equity 100, two additional visits by year-end yield revenue 71, expense 46, profit 25, cash 224 and contract liability 99. The indirect operating bridge 25+99=124 matches 170 received − 46 paid. Complete construction for cash 60 and three visits for 6 next year: revenue 84+15=99, profit 33 and cash/equity 158. Lifetime profit 58=price 170 − construction 100 − visit costs 12. Keeping old progress would recognise 70 for construction; applying 40\% to the entire 170 would wrongly recognise distinct future visits. Missing allocation/progress or enforceability blocks the numeric transition. The one-year-or-less financing expedient under 63 is explicitly elected for this example, not assumed for long prepaid contracts.

\subsubsection{The invoice and the work need different clocks}
\label{sec:mixed-contract-billing}

The builder can earn the same construction and inspection revenue while being
paid on a different schedule. Retain the preceding allocations, costs and
approved amendment, but place all work and payments in calendar 2024. At
inception and after the amendment, the expected interval between each payment
and its related service is at most one year, and the builder elects the
financing expedient. Receivables and contract assets have a separately assessed
zero impairment allowance in this example. These findings matter: AASB~15
requires the financing assessment and AASB~9 treatment of the rights to
consideration; eventual collection alone establishes neither
\citep[paras.~63,105--109]{aasb15dated}.

Suppose the customer initially advances only 30. Opening cash and capital of
100 become cash 130 and a contract liability of 30. In April the builder pays
construction costs of 40 and visit costs of 2, earning the same revenue of
55 as before. Cash falls to 88. The advance covers the first 30 of that work;
the remaining 25 is a contract asset because the right to payment still
depends on a billing milestone. Profit is 13 and equity is 113.

Now the milestone makes a gross invoice of 80 unconditional. The invoice
credits the advance of 30 already received, leaving a receivable of 50. The
contract asset disappears. But the builder has earned only 55 of the 80 now
paid or due, so a contract liability of 25 also appears. Figure
\ref{fig:mixed-billing-clocks} isolates this change. The receivable and contract
liability coexist: one is an unconditional payment right and the other is
work still owed under the same contract.

\begin{figure}[htbp]
\centering
\begin{tikzpicture}[>=Latex, every node/.style={font=\small},
  stage/.style={draw,rounded corners,align=center,text width=3.25cm,
  minimum height=1.65cm,inner sep=5pt}]
\node[stage] (advance) {Advance 30\\Cash 130\\Contract liability 30};
\node[stage,right=6mm of advance] (work) {Earn 55; pay costs 42\\Cash 88\\Contract asset 25};
\node[stage,right=6mm of work] (invoice) {Invoice 80; credit 30\\Receivable 50\\Contract liability 25};
\draw[->] (advance)--(work);
\draw[->] (work)--(invoice);
\end{tikzpicture}
\caption{The invoice changes the payment right without new work or new cash.
Cash stays at 88 and cumulative revenue at 55 when the final box is reached.}
\label{fig:mixed-billing-clocks}
\end{figure}

The advance must therefore retain its identity after receipt. Crediting it
reduces the new amount collectible; it does not refund cash or recognise
revenue again. Let \(G\) be cumulative gross invoices, \(K\) the advances
credited against those invoices, and \(C\) later collections of the resulting
receivables. The amount still collectible is
\[
  \mathrm{AR}=G-K-C.
\]
For the invoice just issued, this is \(80-30-0=50\). Collecting 20 changes
cash to 108 and the receivable to 30, with no change in revenue. Each credit
names both its advance and invoice; each collection names its invoice. A
second credit cannot reuse an advance that has already been consumed.

To determine the contract balance, count consideration already paid or
unconditionally due once. If \(A\) is cumulative advance cash, adding it to
gross invoices would count credited advances twice. Subtract \(K\) to obtain
\(T=G+A-K\). Compare this with cumulative earned revenue \(R\):
\[
  T=G+A-K,\qquad
  \mathrm{CA}=\max(R-T,0),\qquad
  \mathrm{CL}=\max(T-R,0).
\]
At invoicing, \(T=80+30-30=80\), giving the liability \(80-55=25\).
Collection leaves \(T\) unchanged: cash replaces a right that was already
unconditional. Together these relations imply, for cumulative paid costs
\(E\) and opening capital 100,
\[
  \text{cash}+\mathrm{AR}+\mathrm{CA}-\mathrm{CL}
       =100+R-E.
\]
Substitution shows why: cash is \(100+A+C-E\), the receivable is
\(G-K-C\), and the net contract asset is \(R-G-A+K\). The advance,
invoice, credit and collection terms cancel, leaving earned profit in equity.
The cancellation checks the arithmetic; the contractual findings determine
which amounts belong in it.

The May amendment now adds construction revenue of 6 without a receipt,
reducing the liability from 25 to 19. A new, separately identified advance of
50 raises cash to 158 and the liability to 69. Two further visits earn 10
and cost 4, leaving cash 154, receivables 30, contract liability 59 and profit
25. The remaining gross invoice of 90 credits the new advance of 50. Its net
right of 40 increases total receivables to 70; paid-or-due consideration
reaches 170, so the contract liability becomes \(170-71=99\).

Collect the first invoice's remaining 30, then complete construction and the
last three visits for cash costs of 66. Completion earns the remaining 99 and
removes the contract liability. Cash is 118, the second invoice still carries
40, and equity is 158. Collecting that 40 leaves cash and equity of 158, as
in the prepaid example. Customer cash is advances \(30+50\) plus collections
\(20+30+40\), totalling 170. Paid costs are 112 and profit is 58. The gross
invoices of 170 describe the rights created, not an additional receipt.

Full prepayment provides a useful limiting case. If the final invoice credits
all 170 already received, it creates no receivable and no monetary journal;
it still records which advances the invoice consumes. Conversely, invoicing
before any work can create a receivable and contract liability together.
Both schedules preserve the same earned amounts under these terms. The
finite method retains these billing identities alongside the two performance
calculations. Nonzero impairment, refundable excess payments, contract losses
and tax require their own measurements before this example can be joined to
a more general company book.


\hypertarget{unfinished-production-carries-component-specific-completion}{%
\subsubsection{Unfinished production carries component-specific completion}\label{unfinished-production-carries-component-specific-completion}}

The production record adds started/completed/scrapped quantities, material insertion point, conversion stages, opening WIP by component, observed progress, actual eligible costs, normal capacity, fixed/variable overhead classification, NRV, and the costing-policy version. Equivalent units are a derived policy justified by this homogeneous process; IAS~2 \citep[IAS~2.12--18]{eu20231803} prescribes rational cost allocation and the normal-capacity constraint, not this exact spreadsheet convention.

Start 100 units with no opening WIP or abnormal waste. Finish 80;20 are half converted, with all material already inserted. Material 1000 is spread over 100 material units; variable conversion 450 over 90 equivalent conversion units. Fixed cash overhead 600 has admitted normal capacity 120 equivalent units, so its rate is 5, assigned 450 and expensed idle 150. The material denominator is 100 units, while conversion uses only 90 equivalent units:
\[80+20\left(\frac12\right)=90,\qquad\text{completed unit cost}=10+5+5=20.\]
Separate denominators keep missing conversion work out of inventory cost. Completed cost 1600 and WIP 300 exhaust assigned cost 1900: WIP has material 200 and conversion 100. Sell 50 finished units at 30: revenue 1500, cost 1000 and idle expense 150 give profit 350; finished 600 plus WIP 300 remain. Opening cash/equity 3000 leaves cash 2450 and equity 3350.

In the next production month, the remaining 10 equivalent conversion units require variable 50 and fixed 50; independently supplied normal capacity 10 admits both. Add 100 to WIP 300 and transfer completed 400. The remaining 50 finished units then sell 1500 at total cost 1000. Cumulative cash/equity 3850 and profit 850 follow; no WIP remains. A different opening WIP, spoilage policy or completion pattern needs its own allocation. Reusing 100 as conversion units overstates completion; capitalising all 600 fixed cost conceals 150 idle expense. NRV is separately assessed for both unfinished and completed stock, including costs to complete.

\hypertarget{joint-products-and-delayed-supplier-payment}{%
\subsubsection{Joint products and delayed supplier payment}\label{joint-products-and-delayed-supplier-payment}}

The joint-production record retains inseparable joint costs, split-off quantities, separate subsequent costs, allocation-basis evidence, selling values/dates and NRV. For joint cost 300, A 60 at split-off price 4 and B 40 at 9 give values 240/360. The selected rational and consistent relative-sales method assigns 120/180, or unit costs 2/4.5. Selling 30 A and 10 B gives revenue 210, cost 105 and profit 105. Opening cash 500 leaves cash 410, unsold A 60/B 135 and equity 605. Weighting by mass instead changes the admitted policy; equal weights are not a missing-value fix.

Separately, production 300 with an independently established immaterial byproduct of NRV 20 allocates main inventory 280 and byproduct 20. Sale of that byproduct for 20 removes its 20 carrying amount; it earns no margin here. Opening cash 500 ends cash 220+main 280=equity 500. Retaining main 300 as well as byproduct 20 invents cost; deducting 20 without recording the unsold byproduct loses an asset.

The supplier record names the original creditor, normal-credit purchase price, actual terms, financing assessment, initial liability, effective rate, payment schedule and actual payments. Normal cash price 100 and a fixed payment 121 in two years imply 10\% effective interest under the supplied terms. Record inventory and payable 100; interest 10 then 11 raises the liability 121. Ordinary inventory is not a qualifying asset in this case. Sell for cash 150 in year 1: opening cash/equity 100 becomes cash 250, payable 110 and equity 140. Year 2 interest 11 and payment 121 leave cash/equity 129; lifetime profit 29. Under the explicit historical IAS 7 election, both original supplier payment and interest are operating cash. Year 1 indirect 40+payable increase 110=150; year 2 − 11 − 110= − 121. A substituted finance provider or altered arrangement requires reclassification analysis.

\hypertarget{uncertain-tax-has-its-own-measurement-decision}{%
\subsubsection{Uncertain tax has its own measurement decision}\label{uncertain-tax-has-its-own-measurement-decision}}

The prior tax dossier already specifies tax bases, gross DTA/DTL, recoverability, rate changes and lease interaction; the error/estimate bridge retains equity origins. Uncertain tax treatment was genuinely deferred. Under the historical EU IFRIC~23.6--14 \citep[IFRIC~23.6--14,A1--A5]{eu20231803} method, the uncertain-tax record adds treatment grouping, filing position, authority/examination scope, full-information acceptance finding, candidate outcomes, predictive-measure choice, consistent current/deferred bases, reassessment event and disclosure obligations. Assume examination with full knowledge; do not multiply exposure by an estimated chance of avoiding examination.

Synthetic law ultimately deducts equipment 100, with the timing uncertain. Book depreciation 20 leaves PPE 80; cash operating revenue 100 makes pretax profit 80. Opening cash/equity 500, purchase 100 and cash revenue 100 leave cash 500. The enacted rate 25\% is supplied, all tax amounts remain payable, and all origins are profit. If full deduction 100 is probably accepted, current tax 0 and DTL 20 follow. If outcomes 100/20 are binary with full acceptance probability 1/5, and most-likely amount better predicts resolution, use deduction 20: current 20 and DTL 0. In a third branch, deductions 100/40/20 at 1/5,3/10,1/2 and a supported expected-value choice give an expected deduction of

\[\frac15(100)+\frac3{10}(40)+\frac12(20)=42.\]

Taxable profit and the remaining tax base are both 58. Current tax is \(58(25\%)=14.5\), while the taxable temporary difference is \(80-58=22\), giving a deferred tax liability of \(22(25\%)=5.5\). All branches have total tax 20, profit 60 and equity 560; they differ in the timing liability. These equal totals arise from this particular full-deduction law.

If new evidence immediately resolves the expected-value branch at deduction 40, increase current tax 14.5 → 15 and reduce DTL 5.5 → 5. The two 0.5 expense/benefit entries net to zero. Pay 15: cash 485+PPE 80=DTL 5+equity 560. A later resolution must also include intervening depreciation and other facts; silence alone is not new evidence of acceptance. Real law, probability and predictive suitability remain required. The printed hypothetical rate or distribution cannot fill them.

\hypertarget{combining-the-new-operating-contributions}{%
\subsubsection{Combining the new operating contributions}\label{combining-the-new-operating-contributions}}

Combine only the five explicitly compatible EU year 1 contributions: WIP, joint products, byproduct, supplier financing and lifetime restoration. Give each independent lots/contracts/assets but one reporting entity, currency and period; sum their openings, without intercompany elimination or transferring AU cases. Opening cash 4200, PPE 120, provision 20 and equity 4300 become cash 3430, inventory 1375, PPE 71, liabilities 143 and pretax equity 4733. Pretax profit 433 reconciles equity 4300 → 4733; total assets 4876 equal liabilities 143+equity 4733. The same cash movement follows from the indirect reconciliation:
\[433+60-1375+110+2=-770.\]
Starting from pretax profit 433, add depreciation 60, deduct inventory growth 1375, add the supplier liability increase 110 and add noncash restoration interest 2. The 11 restoration estimate adds both PPE and provision and is excluded from cash. This is a pretax operating slice; new tax effects are unresolved, so a complete after-tax report cannot be issued from it. The separate uncertain-tax example is not silently grafted onto this different law/fact pattern.
